\chapter{Conception}

\section{Introduction}
La phase de conception constitue le pont entre l'analyse des besoins et l'implémentation technique. Elle définit l'architecture logicielle, le modèle de données NoSQL et les interactions dynamiques du système eduGenius.

\section{Architecture Logicielle}
\subsection{Architecture Globale (N-Tier)}
Le système repose sur une architecture à trois niveaux :
\begin{enumerate}
    \item \textbf{Niveau Présentation (Frontend)} : Développé avec \textbf{Next.js 16}, il assure une interface utilisateur fluide et réactive.
    \item \textbf{Niveau Application (Backend)} : Un serveur \textbf{Node.js/Express} qui expose une API RESTful et gère les WebSockets via Socket.io.
    \item \textbf{Niveau Données} : \textbf{MongoDB Atlas} pour la persistence des données et \textbf{AWS S3} pour le stockage des objets volumineux.
\end{enumerate}
\section{Conception des Logiques Métiers Complexes}
\subsection{Calcul Automatique des Moyennes et PV}
L'un des modules les plus complexes d'eduGenius est la génération du \textbf{Procès-Verbal (PV)}. La logique de calcul suit les standards du système LMD tunisien.

\begin{table}[H]
\centering
\begin{tabular}{|l|p{10cm}|}
\hline
\textbf{Composant} & \textbf{Formule / Règle de Calcul} \\ \hline
Note de Matière & $(Note_{Examen} \times 0.7) + (Note_{CC} \times 0.3)$ \\ \hline
Moyenne Semestre & $\sum (Note_{Matière} \times Coeff) / \sum Coeff$ \\ \hline
Mention & Passable ($\geq 10$), Assez Bien ($\geq 12$), Bien ($\geq 14$), Très Bien ($\geq 16$), Excellent ($\geq 18$) \\ \hline
Statut Session & Admis (ADM), Ajourné (AJ), Session Spéciale \\ \hline
\end{tabular}
\caption{Règles de gestion pour la génération automatique du PV}
\end{table}

\subsection{Modèle de Gamification (XP System)}
Pour stimuler l'engagement, nous avons implémenté un moteur de récompense basé sur les points d'expérience (XP).
\begin{itemize}
    \item \textbf{Quiz Complété} : $+100$ XP $\times$ Score (\%).
    \item \textbf{Résumé Consulté} : $+20$ XP.
    \item \textbf{Assiduité Parfaite (Semaine)} : $+500$ XP de bonus.
\end{itemize}

\section{Conception de la Base de Données}
...
\subsubsection{Collection : AcademicYears}
Cette collection définit les fenêtres temporelles du système.
\begin{table}[H]
\centering
\begin{tabular}{|l|l|p{8cm}|}
\hline
\textbf{Champ} & \textbf{Type} & \textbf{Description} \\ \hline
year & String & Libellé de l'année (ex: 2023-2024). \\ \hline
isCurrent & Boolean & Définit l'année active par défaut. \\ \hline
semesters & Array & Tableau d'objets (Dates de début/fin, numéro). \\ \hline
\end{tabular}
\caption{Dictionnaire de données - Collection AcademicYears}
\end{table}

\subsection{Dictionnaire de Données NoSQL Détaillé}

\subsubsection{Choix du Modèle : Normalisation vs Dénormalisation}
Dans eduGenius, nous avons privilégié la \textbf{dénormalisation} pour les données consultées fréquemment (ex: inclure le nom du professeur dans le document Cours) afin de minimiser les jointures complexes (lookups) dans MongoDB. En revanche, pour les données académiques sensibles (Notes, Présences), nous utilisons des références d'ID (\textbf{Normalisation}) pour garantir l'intégrité et la cohérence.

\subsubsection{Collection : Departments}
\begin{table}[H]
\centering
\begin{tabular}{|l|l|p{8cm}|}
\hline
\textbf{Champ} & \textbf{Type} & \textbf{Description} \\ \hline
name & String & Libellé du département (ex: Génie Informatique). \\ \hline
head & ObjectId & Référence vers l'enseignant chef de département. \\ \hline
subjects & Array & Liste des matières rattachées. \\ \hline
\end{tabular}
\caption{Dictionnaire de données - Collection Departments}
\end{table}

\subsubsection{Collection : Subjects (Matières)}
\begin{table}[H]
\centering
\begin{tabular}{|l|l|p{8cm}|}
\hline
\textbf{Champ} & \textbf{Type} & \textbf{Description} \\ \hline
name & String & Nom de la matière (ex: Algorithmique). \\ \hline
coefficient & Number & Poids dans la moyenne (Système LMD). \\ \hline
credits & Number & Nombre d'unités de valeur ECTS. \\ \hline
\end{tabular}
\caption{Dictionnaire de données - Collection Subjects}
\end{table}

\subsubsection{Collection : Announcements}
\begin{table}[H]
\centering
\begin{tabular}{|l|l|p{8cm}|}
\hline
\textbf{Champ} & \textbf{Type} & \textbf{Description} \\ \hline
title & String & Titre de l'annonce. \\ \hline
content & String & Corps du message. \\ \hline
isPinned & Boolean & Définit si l'annonce est épinglée en haut. \\ \hline
classId & ObjectId & Classe cible de l'annonce. \\ \hline
\end{tabular}
\caption{Dictionnaire de données - Collection Announcements}
\end{table}

\subsubsection{Collection : Attendance (Présences)}
\begin{table}[H]
\centering
\begin{tabular}{|l|l|p{8cm}|}
\hline
\textbf{Champ} & \textbf{Type} & \textbf{Description} \\ \hline
studentId & ObjectId & Référence vers l'étudiant. \\ \hline
sessionId & ObjectId & Référence vers la session vidéo. \\ \hline
durationInMinutes & Number & Temps total passé dans la salle. \\ \hline
status & Enum & Present, Absent, Late. \\ \hline
\end{tabular}
\caption{Dictionnaire de données - Collection Attendance}
\end{table}
\section{Conception de l'Infrastructure Cloud}
L'infrastructure d'eduGenius a été conçue pour être résiliente et élastique. Nous utilisons une approche \textbf{Cloud Native}.

\subsection{Schéma de Déploiement (Cloud Architecture)}
\begin{itemize}
    \item \textbf{VPC (Virtual Private Cloud)} : Pour isoler les ressources backend et la base de données.
    \item \textbf{Load Balancer} : Répartition de la charge entre plusieurs instances du serveur Express en cas de pic de trafic.
    \item \textbf{CDN (Content Delivery Network)} : Distribution globale des assets statiques (Frontend Next.js) via Vercel.
\end{itemize}

\section{Architecture de Sécurité}
\subsection{Sécurisation des Données en Transit}
Toutes les communications entre le client et le serveur, ainsi qu'entre le serveur et les API tierces (AWS, Gemini), sont chiffrées via le protocole \textbf{TLS 1.3} (HTTPS).

\subsection{Sécurisation des Données au Repos}
\begin{itemize}
    \item \textbf{MongoDB Atlas} : Chiffrement automatique des disques (\textit{Encryption at Rest}).
    \item \textbf{AWS S3} : Utilisation de politiques de bucket (\textit{IAM Policies}) restrictives pour empêcher l'accès public aux documents de cours.
    \item \textbf{Hachage des Mots de Passe} : Utilisation de l'algorithme \textbf{Argon2} ou \textbf{bcrypt} avec un sel (salt) dynamique pour prévenir les attaques par force brute.
\end{itemize}

\section{Architecture des Micro-Services d'IA}
\subsection{Conception du Moteur d'IA Souverain}
Contrairement aux architectures Cloud, eduGenius repose sur un modèle \textbf{Side-car}. Le serveur Node.js communique en local (localhost) avec l'instance \textbf{Ollama}.

\subsubsection{Workflow de Traitement Local}
1. \textbf{Réception} : Le contrôleur reçoit la demande (Résumé/Quiz).
2. \textbf{Extraction} : Le texte brut est extrait via \texttt{pdf-parse}.
3. \textbf{Inférence} : Une requête HTTP POST est envoyée à l'API locale d'Ollama sur le port 11434.
4. \textbf{Streaming} : Le flux de réponse de Gemma 3 est capturé par le backend et stocké après validation JSON.

\fbox{Placeholder: Schéma de l'Architecture Local AI (Node.js <-> Ollama Gemma)}
Cette architecture garantit que même sans accès externe, les enseignants de Polytech-Intl peuvent continuer à générer des quiz instantanément.

\section{Conception des Interactions (Diagrammes de Séquences)}
\subsection{Génération de Résumé IA}
1. Client -> POST /summary.
2. Server -> Fetch PDF from S3.
3. Server -> AI Service (Gemini).
4. Server -> Save result to MongoDB.
5. Server -> Response 200 OK.
\fbox{Placeholder: Diagramme de Séquence - Génération de Résumé}

\section{Conclusion}
Ce chapitre a détaillé les choix structurants d'eduGenius. En passant d'une architecture monolithique à une approche orientée services et en exploitant la flexibilité du NoSQL, nous avons bâti un système capable de supporter les exigences du temps réel et de l'intelligence artificielle. La conception rigoureuse des schémas et des flux d'interaction garantit que la plateforme reste performante même sous une charge importante. Le chapitre suivant portera sur la mise en œuvre concrète de ces concepts et la présentation des résultats obtenus.
